\DocumentMetadata{
  pdfversion=2.0,
  pdfstandard=ua-2,
  lang=en-US,
  tagging=on,
  tagging-setup={math/setup=mathml-SE,tabsorder=structure}
}
\documentclass[11pt,letterpaper]{article}
\usepackage[margin=0.68in,includefoot]{geometry}
\usepackage{newcomputermodern}
\usepackage{amsmath,amscd,mathtools}
\usepackage{tikz,adjustbox}
\providecommand{\square}{\mdlgwhtsquare}
\usepackage[hidelinks]{hyperref}
\hypersetup{
  pdftitle={Analysis First-Year Exam, January 2014, Part 2},
  pdfauthor={University of Florida Department of Mathematics},
  pdfsubject={Graduate mathematics examination},
  pdfdisplaydoctitle=true,
  bookmarksopen=true
}
\setcounter{secnumdepth}{0}
\setlength{\parindent}{0pt}
\setlength{\parskip}{0.28em}
\setlength{\emergencystretch}{3em}
\setlength{\leftmargini}{2.15em}
\setlength{\leftmarginii}{2.65em}
\setlength{\leftmarginiii}{3em}
\renewcommand{\labelenumi}{\textbf{\theenumi.}}
\pagestyle{empty}
\raggedbottom
\allowdisplaybreaks
\newenvironment{examproblems}[1]{%
  \begin{enumerate}%
  \setcounter{enumi}{#1}%
  \setlength{\topsep}{0.35em}%
  \setlength{\partopsep}{0pt}%
  \setlength{\itemsep}{0.48em}%
  \setlength{\parsep}{0.18em}%
}{\end{enumerate}}
\newenvironment{examsubparts}{%
  \begin{enumerate}%
  \setlength{\topsep}{0.18em}%
  \setlength{\partopsep}{0pt}%
  \setlength{\itemsep}{0.16em}%
  \setlength{\parsep}{0.1em}%
}{\end{enumerate}}
\newenvironment{examinstructions}{%
  \par\begingroup\itshape
}{%
  \par\endgroup\vspace{0.18em}
}
\makeatletter
\renewcommand\section{\@startsection{section}{1}{0pt}%
  {-1.25ex plus -.25ex minus -.1ex}%
  {0.55ex plus .1ex}%
  {\normalfont\large\bfseries}}
\renewcommand{\@maketitle}{%
  \newpage\null\vskip -1em%
  {\footnotesize\bfseries
    \href{https://www.ufl.edu/}{University of Florida}, \href{https://math.ufl.edu/}{Department of Mathematics}\par}%
  \vskip -0.5em%
  \tagstructbegin{tag=Title}%
    {\LARGE\bfseries\href{https://gma.math.ufl.edu/exams/analysis-first-year/}{Analysis First-Year Exam}, January 2014, Part 2\par}%
  \tagstructend%
  \vskip 0.25em%
  \tagmcbegin{artifact}%
    \hrule height 1pt%
  \tagmcend%
  \vskip 1em%
}
\makeatother
\title{Analysis First-Year Exam, January 2014, Part 2}
\author{}
\date{}
\begin{document}
\maketitle
\thispagestyle{empty}
\begin{examinstructions}
Answer FOUR questions. Write in a neat and logical fashion, giving complete reasons for all steps.
\end{examinstructions}
\begin{examproblems}{0}
\item
Let~\(f\) be a continuous real-valued function on \([0,1]\). For each \(n \geq 1\), let \(h_n=n\mathbf{1}_{[0,1/n]}\) be~\(n\) times the indicator function of \([0,1/n]\). Prove that 
\[
\lim_{n\to\infty}\int_0^1 f(t)h_n(t)\,dt=f(0).
\]
\item
Let~\(f\) be a continuous real-valued function on \([0,1]\). For \(n\geq 1\) and \(0\leq t\leq 1\), let \(f_n(t)=f(t)t^n\). Determine a condition on \(f(1)\) that is necessary and sufficient for the sequence \((f_n)_{n=1}^{\infty}\) to be uniformly convergent; be sure to prove necessity and sufficiency.
\item
For \(n\geq 1\), let \(f_n\) be a continuous real-valued function on \([0,1]\), and for \(0\leq t\leq 1\) define 
\[
F_n(t)=\int_0^t f_n(u)\,du.
\]
 In each of the following cases, decide whether \((F_n)_{n=1}^{\infty}\) must have a uniformly convergent subsequence, giving proof or counterexample as appropriate:
\begin{examsubparts}
\item[\textnormal{(i)}]
\(|f_n|\leq g\) for some continuous~\(g\) and all~\(n\);
\item[\textnormal{(ii)}]
\(0\leq f_1\leq f_2\leq\cdots\).
\end{examsubparts}
\item
Let \((f_n)_{n=1}^{\infty}\) be a sequence of measurable functions on some measure space. Show that each of the following sets is measurable: 
\[
A=\{\omega\mid f_n(\omega)>0\text{ for infinitely many }n\};
\]
 
\[
B=\{\omega\mid f_n(\omega)\leq 0\text{ for finitely many }n\}.
\]
 Warning: \(A\) and~\(B\) need not be complementary!
\item
Let \((f_n)_{n=1}^{\infty}\) be a uniformly-bounded sequence of continuous functions on \([0,1]\) that converges pointwise to zero. Does it follow that 
\[
\lim_{n\to\infty}\int_0^1 f_n(t)\,dt=0?
\]
 Does the answer change if ‘uniformly-bounded’ is removed from the hypotheses? Give proof or counterexample as appropriate.
\end{examproblems}
\end{document}
